\subsection{A nonflat ring with no degree-two obstruction}

\begin{proposition}[Proposed: cancellation across a nondiscrete valuation branch]
\label{C:nondiscrete-torsion-C2}
Let $A$ be a characteristic-zero valuation domain with maximal ideal
$\mathfrak m_A$, residue field $\kappa$ of characteristic two, and value
group having no least positive element. Let $B$ be a valuation domain
of characteristic two, with residue field $\kappa$ and a specified
coefficient-field section $\kappa\to B$. Set
\[
                         R=A\times_\kappa B.
\]
Then projection and its coefficient-field section induce inverse
isomorphisms
\[
                       R_R(C_2)\simeq R_A(C_2).
\]
Consequently the degree-two polynomial map for $R$ is an isomorphism.
If $B\ne\kappa$, the ring $R$ is reduced, local and of characteristic
zero, but is not flat over $\mathbb Z$: every nonzero element of
$0\times\mathfrak m_B$ is killed by two.
No claim about all symmetric-group degrees is made.
\end{proposition}
\begin{proof}
The maps $R\to A$ and
$a\mapsto(a,\overline a\cdot1_B)$ are split ring maps. Put
$J=0\times\mathfrak m_B$ and regard $A$ as a subring of $R$ by this
section. Then $R=A\oplus J$ additively and $\mathfrak m_AJ=0$.
The ring $R$ is local with residue field $\kappa$, so all finite
projective underlying modules are free.

Let $M=(R^r,T)$ be a $C_2$-lattice. First remove all regular summands
of its residual representation. To justify this directly, if a vector
$\bar v$ belongs to a residual Jordan block of size two, lift it to
$v\in M$. The vectors $v,Tv$ are linearly independent modulo the
maximal ideal, so their span is a free direct summand of the underlying
module and is equivariantly $R[C_2]$. Choose an $R$-linear functional
$\lambda:M\to R$ with $\lambda(v)=1$ and $\lambda(Tv)=0$.
Then
\[
             m\longmapsto\lambda(m)v+\lambda(Tm)Tv
\]
is an equivariant retraction onto this regular summand. Repeat until
$\overline T=I$. Regular summands are already extended from $A$.

Write $T=T_0+U$, where $T_0$ is the matrix over $A$ obtained by
projection and every entry of $U$ lies in $J$. We have $T_0^2=I$
and $T_0\equiv I\pmod{\mathfrak m_A}$.
Since only finitely many entries occur and the value group has no
least positive element, choose $0\ne c\in\mathfrak m_A$ whose
valuation is strictly smaller than the valuation of every nonzero
entry of $T_0-I$ and also strictly smaller than $v(2)$. Then
\[
           D=(T_0-I)/c\in M_r(\mathfrak m_A),
           \qquad DU=UD=0,\qquad cU=0.
\]
The strict inequality for $v(2)$ is harmless and makes clear that the
construction uses the arbitrarily small valuation scale; the displayed
integrality of $D$ is the condition used below.

Consider the block action and change of basis
\[
 E=\begin{pmatrix}T&cI_r\\0&-T_0\end{pmatrix},
 \qquad P=\begin{pmatrix}I_r&0\\-D&I_r\end{pmatrix}.
\]
Both diagonal blocks of $E$ square to one, and its off-diagonal square
is $c(T-T_0)=0$. Thus $E$ is a genuine involution and gives an
admissible sequence with endpoint actions $T,-T_0$.
Using $cD=T_0-I$, $T_0D=-D$ and $DU=UD=0$, direct multiplication gives
\[
                 P^{-1}EP=
                  \begin{pmatrix}I_r+U&cI_r\\0&-I_r\end{pmatrix}.
\]
This is a second admissible sequence. Therefore
\[
                  [M]+[-T_0]=[I_r+U]+r[\mathrm{sgn}].       \tag{C.NT1}
\]

We claim $[I_r+U]=r[\mathbf1]$. On the $A$ branch this action is
trivial. On the $B$ branch it is an involution, hence it has the form
$I+N$ with $N^2=0$. Over a valuation domain every finitely generated
torsion-free module is free; therefore $\operatorname{im}N$ is finite
free, and the exact sequence
\[
       0\longrightarrow\ker N\longrightarrow B^r
          \longrightarrow\operatorname{im}N\longrightarrow0
\]
splits over $B$. In particular $\ker N$ is finite free. Choose a
$B$-basis beginning with a basis of $\ker N$. The matrix of $N$ is
strictly upper triangular (its only possibly nonzero block maps the
chosen complement into the kernel). Lift the residual change-of-basis
matrix to an invertible matrix over $A$; together the two matrices
form a change of basis over $R$. Its $A$ action remains the identity,
and its $B$ action is upper triangular with all diagonal entries one.
Thus the resulting $R$ action has a full admissible flag of trivial
free lines, proving the claim.

Finally the standard regular-tensor exact sequence gives
\[
             [T_0]+[-T_0]=r([\mathbf1]+[\mathrm{sgn}])
\]
over $A$, and hence over $R$. Substitute this and the claim into
(C.NT1) to obtain $[M]=[T_0]$. Together with the previously removed
regular summands, this proves that projection and section are inverse
on every class. The valuation-domain criterion
\ref{C:valuation-domain-criterion} gives the degree-two polynomial
assertion over $A$, and hence over $R$.

The ring $R$ is a subring of the product of two domains and is
therefore reduced. Its projection onto $A$ and the given section
show that it has characteristic zero. If $0\ne b\in\mathfrak m_B$,
then $(0,b)\ne0$ and $2(0,b)=0$, proving failure of flatness over
$\mathbb Z$.
\end{proof}

An explicit example is obtained by choosing compatible roots
$\pi_0=2$, $\pi_{n+1}^2=\pi_n$ in an algebraic closure of $\mathbb Q_2$
and setting
\[
 A=\bigcup_{n\ge0}\mathbb Z_2[\pi_n],\qquad
 B=\mathbb F_2[x]_{(x)},\qquad R=A\times_{\mathbb F_2}B.
\]
The rings $\mathbb Z_2[\pi_n]$ are DVRs with uniformizer $\pi_n$,
and their compatible union is a valuation domain with value group
$\mathbb Z[1/2]$ and residue field $\mathbb F_2$. Hence the proposition
applies, while $(0,x)$ is a nonzero element killed by two. This is an
all-rank degree-two statement only, not a positive solution of C.1
for this coefficient ring.
